% =====================================================================
%  基于用户评分的生成智能体后训练 —— SFT / 奖励模型 / 偏好优化 方法手册
%
%  编译(必须 XeLaTeX 或 LuaLaTeX,含中文):
%      xelatex rm_posttraining.tex && xelatex rm_posttraining.tex
%  依赖:ctex, amsmath, amssymb, booktabs, geometry(TeX Live 标配)。
%
%  内容:以"剧本创作智能体 + 用户满意度评分"为场景,系统整理
%  (1) 评分信号的性质与整形;(2) SFT 两种用法;(3) 奖励模型的标准
%  训练法(结构/损失/日程/验收/校准);(4) 偏好优化三路线
%  (DPO / KTO / RM+RL)及流程;(5) 风险防线与选择判据。
%  公式均带编号与标签,可整段摘抄。
% =====================================================================
\documentclass[UTF8, 11pt]{ctexart}

\usepackage{amsmath}
\usepackage{amssymb}
\usepackage{booktabs}
\usepackage[a4paper, margin=2.5cm]{geometry}

\DeclareMathOperator{\logsigmoid}{log\,\sigma}
\allowdisplaybreaks

\title{基于用户评分的生成智能体后训练\\[4pt]
       \large SFT · 奖励模型 · 偏好优化 —— 方法手册}
\author{}
\date{}

\begin{document}
\maketitle
\vspace{-2em}
\tableofcontents
\newpage

% =====================================================================
\section{问题设定与信号性质}
% =====================================================================

\subsection{设定}

生成智能体(以剧本创作为例)接收提示 $x$(用户想法),输出文本 $y$(剧本);
线上收集到用户满意度评分 $s\in\{1,\dots,5\}$。目标:利用三元组数据
$\mathcal{D}=\{(x_i,\,y_i,\,s_i)\}$ 通过后训练提升生成质量。

\subsection{评分信号的四个性质(方法设计的出发点)}

\begin{enumerate}
  \item \textbf{稀疏}:一条产出一个标量,数据量远低于监督训练所需;
  \item \textbf{尺度不一}:同一分值在不同用户处含义不同(评分习惯差异);
  \item \textbf{噪声大}:主观口味,同人同稿不同时刻打分亦漂移;
  \item \textbf{有偏}:愿意评分的用户并非随机样本。
\end{enumerate}

由此得出贯穿全文的原则:\textbf{人类评分的``序''可靠、``值''不可靠;
一切设计围绕多榨序、少信值。}

\subsection{数据整形}

\paragraph{用户内标准化.} 对用户 $u$ 的全部评分做 z-score:
\begin{equation}
\tilde s \;=\; \frac{s-\mu_u}{\sigma_u+\varepsilon},
\label{eq:zscore}
\end{equation}
消除逐用户的尺度平移与伸缩;评分人数过少或方差异常的样本剔除。

\paragraph{三种目标形态.} 整形后的数据组织为:
(i) \textbf{高分子集} $\mathcal{D}^{+}=\{(x,y)\,\vert\,\tilde s\ge\tau\}$;
(ii) \textbf{偏好对} $\{(x,\,y_w,\,y_l)\}$ —— 同一提示下两份产出,
$\tilde s_w-\tilde s_l\ge\delta$ 方成对(分差过小视为噪声);
(iii) \textbf{二元标签} $\{(x,\,y,\,z)\},\ z\in\{\text{好},\text{差}\}$ ——
最常见的真实形态(一条提示只有一份产出,凑不成对)。
三种形态分别对应第 2、4.1、4.2 节的方法。

% =====================================================================
\section{SFT 路线}
% =====================================================================

\subsection{拒绝采样 SFT(Filtered / Rejection-sampling SFT)}

只在高分子集上做极大似然:
\begin{equation}
\mathcal{L}_{\text{SFT}}(\theta)
=-\sum_{(x,y)\in\mathcal{D}^{+}}\;\sum_{l=1}^{\lvert y\rvert}
\log \pi_\theta\bigl(y_l \mid x,\, y_{<l}\bigr),
\label{eq:sft}
\end{equation}
损失只计生成侧 token(提示部分掩蔽)。要点:阈值 $\tau$ 按\textbf{题材分层}
取分位数而非全局一刀切(防高分题材垄断导致风格坍缩);混入
10--20\% 通用指令数据防灾难性遗忘;LoRA 即可,1--2 个 epoch。

\subsection{条件化 SFT(score-conditioned)}

不丢弃低分样本,把分数作为控制前缀一同训练:
\begin{equation}
\max_\theta\;\sum_{(x,y,s)}\log \pi_\theta\bigl(y \mid [\,\text{分数}{:}\,s\,],\, x\bigr),
\qquad \text{推理时固定条件于最高分。}
\label{eq:condsft}
\end{equation}
低分样本教会模型``什么是差'',全量数据均被利用,数据稀缺时优于过滤式。

% =====================================================================
\section{奖励模型(RM)的标准训练法}
% =====================================================================

\subsection{结构:SFT 底座 + 标量头}

RM 由 SFT 模型改造:去掉词表头(lm\_head),换一个
$\mathbb{R}^{d}\!\to\!\mathbb{R}$ 的线性头;输入 $(x,y)$ 全文,
取\textbf{最后一个有效 token} 的隐藏向量 $h_{\text{last}}$(因果注意力下
唯有末位读完全文)输出标量:
\begin{equation}
r_\phi(x,y) \;=\; w^\top h_{\text{last}}(x,y).
\label{eq:rmhead}
\end{equation}
主体继承 SFT 权重(已``懂''该领域),新增参数仅 $d$ 个;RM 规模通常
不大于被训策略(判断易于生成)。

\subsection{两种损失}

\paragraph{回归(MSE,不推荐为主).}
\begin{equation}
\mathcal{L}_{\text{MSE}}=\bigl(r_\phi(x,y)-\tilde s\bigr)^2 .
\label{eq:mse}
\end{equation}
即便目标是归一化分数,MSE 仍强制拟合绝对值,残余尺度噪声全被当作信号。

\paragraph{Bradley--Terry 成对(工业主流).}
建模``偏好方获胜''的概率只依赖\textbf{分差}:
\begin{equation}
P(y_w \succ y_l \mid x)\;=\;\sigma\bigl(\Delta\bigr),
\qquad
\Delta = r_\phi(x,y_w)-r_\phi(x,y_l),
\label{eq:bt-model}
\end{equation}
对观测到的人类偏好做极大似然,得
\begin{equation}
\boxed{\;
\mathcal{L}_{\text{BT}}
=-\,\mathbb{E}_{(x,y_w,y_l)}\;\logsigmoid\!\bigl(\Delta\bigr)\;}
\label{eq:bt}
\end{equation}
逐步解读:$\sigma$ 把分差映为胜率($\Delta{=}0\Rightarrow 0.5$;
$\Delta{=}2\Rightarrow 0.88$);$-\log$ 是对已发生事实的罚金
(给事实的概率越低罚越重);其对 $\Delta$ 的梯度
\begin{equation}
\frac{\partial \mathcal{L}_{\text{BT}}}{\partial \Delta}
=-\bigl(1-\sigma(\Delta)\bigr)
\label{eq:btgrad}
\end{equation}
的幅值恰为``分给错误结局的剩余概率''——排反的对猛推、已排对且笃定的对
几乎不动,力气自动花在难对上。损失仅含 $\Delta$,对逐用户整体平移免疫
(\textbf{只学序不学值});实现用数值稳定的 \texttt{logsigmoid},
避免 $\sigma$ 下溢后取对数得 $-\infty$。

\paragraph{带间隔变体(Llama-2).}
利用评分幅度:分差大的对要求更大间隔,
\begin{equation}
\mathcal{L}=-\logsigmoid\bigl(\Delta-m(\tilde s_w-\tilde s_l)\bigr),
\qquad m(\cdot)\ \text{单调不减},
\label{eq:margin}
\end{equation}
把被 BT 丢弃的``值''信息部分找回。

\subsection{训练日程与工程要点}

\begin{enumerate}
  \item \textbf{同一排序的所有对置于同一 batch}:一次 $K$ 路排序拆出
        $\binom{K}{2}$ 个高度相关的对,拆散会对同一提示反复过拟合;
        同批处理则每份产出仅前向一次;
  \item \textbf{只训 1 个 epoch}:RM 过拟合极快,第 2 轮起留出准确率回落;
  \item 学习率取 SFT 的 $1/10$ 量级,早停看留出集;
  \item 长度偏置治理:配对时长度配平,或损失中加长度惩罚。
\end{enumerate}

\subsection{验收:成对判对率}

留出偏好对上,RM 给人类偏好方更高分的比例(随机基线 50\%)。
\textbf{65--75\% 即可投产}:主观任务上人类互评一致率亦仅此量级,
它是标签自身的噪声天花板;显著``超越''反而提示过拟合了标注人群癖好。

\subsection{分数尺度的控制(校准)}

BT 训得的分数是\textbf{序数}:零点自由(整体平移不变),
间隔以对数胜率为单位($\Delta{=}1\Leftrightarrow$ 约 73\% 胜率)。
是否需要校准取决于下游:

\begin{center}
\begin{tabular}{@{}lll@{}}
\toprule
下游用途 & 需要绝对值? & 手段 \\
\midrule
同提示 best-of-$N$ & 否 & 无需处理 \\
GRPO 训练 & 否 & 组内 $(r-\bar r)/\sigma$ 已消零点与尺度 \\
PPO(裸奖励) & 半 & 运行时白化(滑动均值/方差) \\
阈值判定 / 对外展示 & 是 & 见下三层 \\
\bottomrule
\end{tabular}
\end{center}

三层校准:(i) \textbf{锚定集归一}——固定代表性样本集,报
$(r-\mu_{\text{anchor}})/s_{\text{anchor}}$,锚定集须版本化;
(ii) \textbf{Platt 校准}——留出对上拟合 $a,b$ 使 $\sigma(a\Delta+b)$
对齐真实胜率(可靠性曲线验收);(iii) \textbf{单调映射回人类分}
(等张回归),供展示;映射必须单调,序不可扭曲。训练期亦可加
$\lambda(r_w^2+r_l^2)$ 的 L2 项预先钉住零点。
\textbf{跨训练轮次禁止直接比较裸分}(两支未对零的温度计)。

\subsection{运行时白化:PPO 消费裸奖励时的尺度治理}

PPO 的目标近似``奖励 $-\ \beta\cdot$KL'',两项相加减,\textbf{数值尺度
必须匹配};而 BT 奖励的尺度任意——分数挤在 $[-0.1,0.1]$ 时奖励项被
KL 项淹没(策略不动),散在 $[-30,50]$ 时缰绳失效(直奔 reward
hacking)。更关键的是分布随训练\textbf{漂移}:策略变好,RM 均分持续
上行,一次性归一化必然过期,故须``运行时''处理。

白化即标准化:
\begin{equation}
r' \;=\; \frac{r-\mu_t}{\sigma_t+\varepsilon},
\label{eq:whiten}
\end{equation}
其中 $\mu_t,\sigma_t$ 为\textbf{滑动统计},随每个 batch 以指数滑动平均
更新:
\begin{equation}
\mu_t=(1-\alpha)\,\mu_{t-1}+\alpha\,\bar r_{\text{batch}},
\qquad
\sigma^2_t=(1-\alpha)\,\sigma^2_{t-1}
+\alpha\,\overline{(r-\mu_t)^2}_{\text{batch}},
\qquad \alpha\in[0.01,\,0.1].
\label{eq:ema}
\end{equation}
变换后进入 PPO 的不再是无意义的裸分,而是``比近期平均好几个标准差''
的无量纲量;分布漂移被滑动窗口自动跟踪($\alpha$ 控制记性长短:过大
随批间噪声抖动,过小跟不上漂移)。数值感受:近期 $\mu{=}3.2$、
$\sigma{=}0.25$ 时,3.5 分折为 $+1.2$;两周后策略进步、$\mu$ 滑至
3.6,同样的 3.5 分折为 $-0.4$——如实反映``此作在当下已属平庸''。

实现细节:常见实现把治理拆在两处——\textbf{奖励}一步只除 $\sigma$、
不减 $\mu$(保留奖励正负语义,防稀疏奖励下``无所事事亦得正分'');
\textbf{优势}另在批内做减均值除标准差,两处合成完整治理。

\paragraph{GRPO 为何天然免此机制.} 组相对优势
$(r^{(i)}-\bar r)/\sigma_r$ 本身就是一次白化,且 $\bar r,\sigma_r$
由\textbf{同一状态的 $K$ 个候选现场算出}——基线精确到``这一题''
而非``最近的平均题'';全局滑动统计所逼近的目标,被组结构一步实现。

% =====================================================================
\section{偏好优化三路线}
% =====================================================================

\subsection{DPO:有成对数据时的首选}

带 KL 约束的 RLHF 目标存在闭式解,最优策略的隐式奖励为
$\beta\log\frac{\pi_\theta}{\pi_{\text{ref}}}$;代回 BT 似然即得
\begin{equation}
\boxed{\;
\mathcal{L}_{\text{DPO}}
=-\,\mathbb{E}_{(x,y_w,y_l)}\log\sigma\!\Bigl(
\beta\Bigl[\log\tfrac{\pi_\theta(y_w\mid x)}{\pi_{\text{ref}}(y_w\mid x)}
-\log\tfrac{\pi_\theta(y_l\mid x)}{\pi_{\text{ref}}(y_l\mid x)}\Bigr]\Bigr)\;}
\label{eq:dpo}
\end{equation}
$\pi_{\text{ref}}$ 为冻结的 SFT 模型;$\beta\in[0.1,0.5]$ 是内建的 KL 缰绳。

\paragraph{偏好对来源(按信号强度).}
(1) 用户编辑对:改后 $\succ$ 原稿(修正行为零噪声);
(2) 重生成行为对:采纳版 $\succ$ 被弃版;
(3) 多稿评分对:同提示多份产出,归一分差超阈方成对;
(4) 主动采集二选一(成本高,但对来自当前策略分布,训练最有效)。

\paragraph{流程.} SFT 复制两份(冻结作 ref / 可训作 $\theta$)
$\to$ 喂偏好对(损失只计生成侧)$\to$ $\beta$ 自 0.1 起调、
lr 低 SFT 一个量级、典型 1 epoch $\to$ 监控隐式奖励差(升)、
chosen 对数似然(不塌)、生成长度(不涨)$\to$ 冻结人评集盲评验收。

\paragraph{坑.} 长度偏置放大(chosen 系统性更长;配平或用长度归一
变体);``双降''现象(两侧 logp 同落,概率质量外漂);多 epoch 必过拟合。

\subsection{KTO:单样本二元标签的正解}

真实数据常凑不成对(一提示一产出)。KTO 依前景理论,
以隐式奖励 $r_\theta=\beta\log\frac{\pi_\theta}{\pi_{\text{ref}}}$
相对参考点 $z_0$(批内 KL 估计)构造:
\begin{equation}
\boxed{\;
\mathcal{L}_{\text{KTO}}
=\mathbb{E}\Bigl[
\lambda_D\bigl(1-\sigma(r_\theta-z_0)\bigr)\Big\vert_{\,z=\text{好}}
+\lambda_U\bigl(1-\sigma(z_0-r_\theta)\bigr)\Big\vert_{\,z=\text{差}}
\Bigr]\;}
\label{eq:kto}
\end{equation}
好样本推上参考点、差样本压下;$\sigma$ 的饱和特性压低极端样本梯度,
对噪声标注稳健。二值化:$\tilde s\ge+0.5\Rightarrow$ 好,
$\tilde s\le-0.5\Rightarrow$ 差,中段丢弃;类别不平衡用
$\lambda_D n_D/(\lambda_U n_U)\in[1,\,4/3]$ 调节而非降采样——
差样本亦被完整利用。

\subsection{RM + 在线 RL:持续爬坡的上限路线}

两阶段:先以第 3 节方法训 RM 并验收;再做在线策略优化。
以 GRPO 为例:同一提示采 $K$ 份产出成组,RM 打分
$\{r^{(i)}\}$,组相对优势
\begin{equation}
A^{(i)}=\frac{r^{(i)}-\bar r}{\sigma_r+\epsilon},
\label{eq:grpoadv}
\end{equation}
以裁剪比值目标更新并以 KL 锚定参考模型(免 critic,契合稀疏标量奖励;
组归一同时消去 RM 分数的零点与尺度,故无需校准)。

\paragraph{Goodhart 双线监控.} RM 分随训练单调上升是常态,
\textbf{以冻结人评集为准}:RM 升而人评平即过优化点,立即停;
防线另有 RM 集成取保守分、奖励中显式长度惩罚、
每轮对新策略产出补标注并\textbf{刷新 RM}(静态 RM 必被钻穿)。

% =====================================================================
\section{扩展:视频生成模型的数据回流}
% =====================================================================

场景变更:不再是文本智能体,而是\textbf{视频生成模型直接出片},由用户
或标注员评分。原则与语言域一脉相承(人类偏好训 RM、RM 驱动优化、
KL 与人评防跑偏),但三个特性要求逐层换零件。

\subsection{与语言域的三点不同}

\begin{enumerate}
  \item \textbf{模型形态}:扩散 / 流匹配模型,\textbf{没有逐 token 似然}
        ——DPO/PPO 依赖的 $\log\pi$ 不可直算,每条后训练路线都须先回答
        ``似然从哪来'';
  \item \textbf{质量多维}:画质、运动质量、文本对齐、时序一致、物理
        合理,单一总分会互相打架(见 5.4 节``静帧 hack'');
  \item \textbf{采样昂贵}:一条样本数十秒至数分钟,在线 RL 的组采样
        成本高一个量级,离线方法(DPO)与拒绝采样性价比更高。
\end{enumerate}

\subsection{视频 RM 的主流训法}

一句话:\textbf{VLM 底座 + 同 prompt 成对标注 + 分维 BT 损失}——
损失仍为式~\eqref{eq:bt},变的是底座与数据。

\begin{enumerate}
  \item \textbf{底座}:视频-语言模型(VLM)加标量头——它已会``看''
        视频并对齐文本,正如语言域的 SFT 底座已会``读''。图像域先例
        ImageReward / PickScore / HPS(CLIP/BLIP 底座 + BT),视频域
        直接承袭;
  \item \textbf{数据协议}:同一 prompt 生成多条,标注员\textbf{按维度
        分别成对比较}——主流如视觉质量(VQ)/ 运动质量(MQ)/
        文本对齐(TA)三维各训一头、加权合成(VideoAlign /
        VideoReward 一系);
  \item \textbf{清单式一派}(VisionReward):不打分,令 VLM 回答数十个
        是/否判断题,线性加权成分——有客观锚的维度用绝对判定;
  \item \textbf{验收}:仍为留出集成对判对率;
  \item \textbf{分工}:标注员产成对、分维的高质量偏好(RM 主粮);
        用户分稀疏混杂,用于筛选高分数据入库与\textbf{校准} RM
        (检验 RM 排序与真实用户口味的一致率)。
\end{enumerate}

\subsection{改进生成模型的四条主流路线(按成本从低到高)}

\begin{center}
\begin{tabular}{@{}p{0.17\textwidth}p{0.2\textwidth}p{0.24\textwidth}p{0.29\textwidth}@{}}
\toprule
路线 & 代表 & 似然问题怎么解 & 要点 \\
\midrule
拒绝采样 SFT & RAFT 式 & 不需要似然 & 每 prompt 生成 $N$ 条,RM 挑
最好者继续微调;最稳最廉,必做的第一步 \\
Diffusion / Flow-DPO & Diffusion-DPO、Flow-DPO & 以 ELBO(去噪损失差)
\textbf{替代}似然 & 偏好对上离线训:令 chosen 的去噪误差降得更多;
图像域自 SD3 一代起标配 \\
可微奖励反传 & ReFL / DRaFT & 不需要似然:RM 可微,梯度\textbf{穿过
去噪链}回传 & 收敛快但 hacking 最凶(过饱和/高对比);须截断反传
步数 + LoRA + 早停 \\
在线 RL & DDPO $\to$ Flow-GRPO $\to$ DanceGRPO & \textbf{ODE$\to$SDE
改写}:每步去噪变为高斯转移,$\log\pi$ 可精确计算 & 同 prompt 采一组、
RM 打分、组相对优势;GRPO 全套机制(比值/裁剪/KL)随即可用 \\
\bottomrule
\end{tabular}
\end{center}

\subsection{落地配方与视频特有防坑}

\begin{enumerate}
  \item 数据:每 prompt 生成 $K$ 条 $\to$ 标注员分维成对比较
        (用户分作校准集);
  \item 训 RM:VLM 底座 + 分维 BT,留出判对率验收;
  \item 先吃便宜的:线上 best-of-$N$ 出片 + 高分产出做拒绝采样 SFT;
  \item 离线一轮 Flow-DPO;
  \item 数据与算力到位后 Flow-GRPO 在线爬坡,KL 锚定参考模型;
  \item 全程:冻结人评集盲评定代际胜负,RM 随策略定期刷新
        (见第 6.2 节)。
\end{enumerate}

\paragraph{视频特有的两类 reward hacking.}
(i) \textbf{静帧 hack}:纯画质 RM 会给``几乎不动''的视频打高分
(无运动即无运动伪影)——这正是运动质量必须单列一维的原因;
(ii) \textbf{过饱和 hack}:可微反传路线的经典病,RM 偏好鲜艳
$\to$ 输出日益``荧光''。防线:多维 RM 互相制衡 + KL + 人评停机线。

% =====================================================================
\section{迭代配方与评估纪律}
% =====================================================================

\subsection{迭代配方}

\begin{enumerate}
  \item 第 1 轮:高分子集 SFT(把分布拉上合格线);
  \item 第 2 轮:线上新评分 $\to$ 构造偏好对 / 二元标签 $\to$ DPO 或 KTO;
  \item 第 3 轮起:部署 $\to$ 收分 $\to$ 再一轮偏好优化(迭代式);
  \item 数据过万且有在线基建后,切换 RM + RL 持续爬坡。
\end{enumerate}

另监控生成长度与熵——长度偏置与风格坍缩的早期信号。

\subsection{两条方向相反的纪律:尺子冻住,判官常换}

\paragraph{冻结人评集盲评(尺子,绝对不变).}
固定一套评测提示词(覆盖各题材),\textbf{永不更改、永不参与训练}。
每出一代模型,新旧两代在同一套提示词上各生成一遍,由人做\textbf{盲评}
A/B:同屏对比、不标来源、左右随机,仅选``哪条更好'';新一代胜率显著
过半方承认进步。三个词各管一件事——\textbf{冻结}管可比性(尺子每代
一换,涨跌说不清是模型变了还是尺子变了);\textbf{盲}管公正(评审知
来源必偏心);\textbf{人评}管真值(线上指标皆可被训练过程污染:RM 分
被策略钻空、用户均分随人群与期望漂移)。铁律:该提示词集绝不可混入
训练数据——进了即泄题,尺子作废。

\paragraph{RM 定期刷新(判官,必须常换).}
RM 训自\textbf{旧策略的产出};RL 一轮之后策略分布迁移——新风格、
新病灶,且策略最擅长的正是寻找 RM 的判卷漏洞(何种色调骗高分、
何种静帧躲扣分)。此时 RM 面对的是分布外样本,打分失真,而策略正沿
失真的高分狂奔——reward hacking 的完整机制即此。刷新即:每轮 RL 后,
对\textbf{当前策略的产出}补一批人类标注(重点覆盖 RM 打高分的区域,
漏洞藏于此),重训或增训 RM,再入下一轮;Llama-2 即以 RM 与策略交替
迭代数轮的方式执行。

\paragraph{为何一冻一换.} 二者职责相反:

\begin{center}
\begin{tabular}{@{}lll@{}}
\toprule
 & 冻结人评集 & 奖励模型 RM \\
\midrule
职责 & 评测:跨代比较 & 奖励:驱动当下优化 \\
需要 & 可比性 $\Rightarrow$ 绝对不变 & 贴近当前分布 $\Rightarrow$ 必须常变 \\
反之则 & 代际胜负失去意义 & 不变才出事:被策略钻穿 \\
\bottomrule
\end{tabular}
\end{center}

一句话:\textbf{尺子一变没法比,判官不换必被骗。}

% =====================================================================
\section{选择判据}
% =====================================================================

\begin{center}
\begin{tabular}{@{}p{0.30\textwidth}p{0.16\textwidth}p{0.44\textwidth}@{}}
\toprule
手头数据形态 & 用什么 & 一句话理由 \\
\midrule
几十条评分 & 上下文回流 & 高分范例进 few-shot,零训练成本 \\
高分子集可辨 & 拒绝采样 SFT & 先把分布拉上合格线 \\
同提示多稿 / 编辑对 & DPO & 有对直训,免 RM,一晚出结果 \\
单稿 + 好/差(最常见) & KTO & 不需成对,差样本不浪费,抗噪 \\
上万条 + 在线基建 & RM + GRPO/PPO & 唯一可持续在线爬坡,须防 Goodhart \\
\bottomrule
\end{tabular}
\end{center}

\vspace{0.5em}
\noindent 贯穿三条铁律:参考模型冻结;损失只计生成侧 token;
验收只认冻结人评集的盲评。

\end{document}
